ここでは、高被引用論文の影響が長期的にどのように変化するかを分析するため、オブソレッセンスに注目し、寿命構造を観察する。
まず、オブソレッセンス状態を判定するために、Sun\cite{SUn1}らの提案する Obsolescence vector を導入する。
Obsolescence vector はGini係数の派生であり$O(Gs,A^-)$で表されるが、本研究では$O(Gs,A^-,A^+)$を用いる。$Gs$は、Gini係数に相当し、値が小さいほど相対的に引用が早期に偏っていることを示し、$|A^-|$、$|A^+|$はローレンツ曲線による面積の一部であり、その値はそれぞれ被引用の減少および増加の変化の速さを示す（図\ref{fig:obsvec}）。
